\documentclass[12pt]{article}
\usepackage[utf8]{inputenc}
\usepackage[T1]{fontenc}
\usepackage[margin=1in]{geometry}
\usepackage{setspace}
\usepackage{xcolor}
\usepackage[colorlinks=true, urlcolor=blue!60!black, linkcolor=blue!60!black]{hyperref}
\setlength{\parindent}{0em}
\setlength{\parskip}{0.6em}
\renewcommand{\baselinestretch}{1.0}
\begin{document}
\begin{center}
{\LARGE \textit{Mahmoud v. Taylor}, 606 U.S. 522 (2025)\par}
\vspace{0.5em}
{\large Case Brief --- Two-Page Edition\par}
\end{center}

\section{Memory Jogger}
Parents won notice and opt-outs from LGBTQ+-inclusive elementary storybooks: a \textit{Yoder}-type burden triggers strict scrutiny without regard to \textit{Smith} (\href{https://www.law.cornell.edu/supremecourt/text/24-297}{Cornell LII}).

\section{Detailed Case Facts}
Montgomery County, Maryland, added five LGBTQ+-inclusive storybooks to its K--5 English curriculum, with guidance to ``[d]isrupt the either/or thinking'' of students. In 2023 it rescinded parental notice and religious opt-outs, while still allowing opt-outs from sex education. Muslim, Catholic, and Orthodox parents sued.

\section{Procedural History}
The district court denied a preliminary injunction; a divided Fourth Circuit affirmed, finding no coercion on a ``threadbare'' record; the Supreme Court reversed.

\section{Judicial Votes}
Reversed, 6--3. Alito wrote for Roberts, Thomas, Gorsuch, Kavanaugh, and Barrett; Thomas concurred; Sotomayor, with Kagan and Jackson, dissented.

\section{Holding}
Denying notice and opt-outs from normative instruction on gender and sexuality substantially interferes with children's religious development under \textit{Yoder}, triggers strict scrutiny regardless of neutrality, and fails narrow tailoring.

\section{Analysis of Opinions}
\textbf{Majority:} For children aged 5 to 11, the ``unmistakably normative'' books posed ``a very real threat of undermining'' parents' teaching; private schooling cannot cure the condition on public education; existing opt-outs showed the burdens were ``self-inflicted.'' \textbf{Thomas:} schools must not get ``a playbook for evading the First Amendment.'' \textbf{Sotomayor:} the Clause bars coercion, not ``exposure''; opt-outs will cause ``chaos.''

\section{Examples: Future Applications}
\textbf{Same-Side:} a second-grade workbook calling Scripture ``just myths,'' with no opt-out, fails. \textbf{Opposite-Side:} a neutral high-school biology unit on evolution survives. \textbf{Fence-Sitter:} a factual seventh-grade health lesson on gender identity.

\section{Critique}
\textit{Mahmoud} narrows \textit{Smith} without overruling it and leaves unclear which burdens are ``\textit{Yoder}-like.'' Separating ``normative'' instruction from ``exposure'' forces judges to parse children's books, and districts may self-censor.

\section{Key Quotations}
\begin{quote}
``When the burden imposed is of the same character as that imposed in \textit{Yoder}, we need not ask whether the law at issue is neutral or generally applicable before proceeding to strict scrutiny. That much is clear from our decisions in \textit{Yoder} and \textit{Smith}.''
\end{quote}
\noindent --- Justice Alito, Opinion of the Court (slip op., at 36).

\end{document}
